\documentclass[11pt,a4paper]{article}
\usepackage[margin=23mm]{geometry}
\usepackage{amsmath,amssymb,amsthm}
\setlength{\emergencystretch}{2em}
\newtheorem{lemma}{Lemma}
\newtheorem{theorem}[lemma]{Theorem}
\newcommand{\gen}[1]{\langle#1\rangle}
\newcommand{\Sym}{\operatorname{Sym}}
\title{Rigidity at the minimum local orbit size}
\author{Human Compression project: L044 continuation}
\date{}
\begin{document}\maketitle
All marked generators are involutions satisfying string
commutations. No intersection property is assumed in
any restricted action below.

\begin{lemma}[The centralizer of a short anchored action]
Let $T=\gen{t_0,\ldots,t_{h-1}}$ be transitive on $V$,
where $h\ge3$, the last mark is nonidentity, and $a$
is fixed by $t_1,\ldots,t_{h-1}$. If $|V|\le h+2$,
the full permutation centralizer of $T$ is trivial.
\end{lemma}
\begin{proof}
The complete acquired anchored-path lemma gives
$|V|\ge h+1$. In degree $h+1$ it identifies
$T=\Sym(V)$, whose centralizer is trivial.
In degree $h+2$ it supplies the distinct path
$a_0=a,a_1,\ldots,a_h$, with $t_i$ exchanging
$a_i,a_{i+1}$ and fixing every earlier path point.
Let $z$ be the remaining point.

Suppose $c\ne1$ centralizes $T$. A centralizer of
a transitive group has no nonidentity element fixing
a point. Thus $c(a)\ne a$. It is fixed by all
$t_1,\ldots,t_{h-1}$. No other path point has this
property: $a_1$ is moved by $t_1$, and $a_k$, $k\ge2$,
is moved by $t_{k-1}$. Consequently $c(a)=z$.
In particular all $t_i$ with $i\ge1$ fix $z$.
The last two marks are now literal overlapping
transpositions $(a_{h-1}\ a_h)$ and
$(a_{h-2}\ a_{h-1})$: every other point is fixed
by their path properties and fixation of $z$.
Their indices are both at least one because $h\ge3$.
Commuting with both transpositions makes $c$
preserve both two-point supports, hence fix their
unique common point $a_{h-1}$. This contradicts
semiregularity and proves the lemma.
\end{proof}

\begin{lemma}[The larger bound with both factors present]
In the complete two-commuting-anchored-factor setting,
assume both active lengths $h,u$ are positive and
$h+u\ge4$. Then the transitive orbit has size at
least $h+u+2$.
\end{lemma}
\begin{proof}
If neither factor is transitive, the acquired count
already gives $2\max(h+1,u+1)\ge h+u+2$.
If the left factor is transitive, all right marks
except the first fix its anchor and centralize it,
so $u=1$. Thus $h\ge3$. The right mark gives a
nontrivial centralizer element, so the preceding
lemma excludes degree at most $h+2$. The degree
is at least $h+3=h+u+2$. Interchange the factors
for the other case. The small pairs of lengths
$(1,1),(1,2),(2,1)$ are outside this statement.
\end{proof}

\begin{theorem}[The minimum orbit and its first omission]
Let $Q=\gen{\rho_m,\ldots,\rho_y}$ be transitive
on $W$, every displayed restriction nonidentity.
Suppose its full local $j$-omission has exactly
two orbits joined by exactly one $j$-edge. Put
$h=j-m$, $u=y-j$ and assume $h,u\ge2$ and
\[
                         |W|=y-m+2=h+u+2.
\]
Then the two local omission orbits $A,B$ can be
named so that $|A|=h+1$, $|B|=u+1$ and
\[
       Q_{<j}=\Sym(A),\qquad Q_{>j}=\Sym(B),
       \qquad Q=\Sym(W),
\]
where each displayed group fixes the complementary
set pointwise. If the first omission $Q_{>m}$
is intransitive, then for one point $z\in A$,
\[
                         Q_{>m}=\Sym(W\setminus\{z\}),
\]
again with $z$ fixed pointwise.
\end{theorem}
\begin{proof}
Use the actual two local omission orbits and their
anchored commuting left and right factors, as in
the complete local orbit-size proof. The greatest
left and right active lengths are $h,u$. They
cannot occur only on the same side: the preceding
lemma would make that side have at least $h+u+2$
points, leaving no point for the other side.
Thus the maxima occur on opposite sides. The
anchored bounds are saturated on both:
$|A|=h+1$ and $|B|=u+1$. The exhausted anchored
path lemma identifies the left image on $A$ as
natural symmetric, and the right image on $B$
as natural symmetric. Both degrees are at least
three. The opposite commuting factor therefore
acts trivially on each side, proving the first
two pure subgroup formulas.

Write the unique crossing as $\{a_0,b_0\}$.
The permutation $\rho_j$ consists of this
transposition and permutations internal to $A$
and $B$. Those internal permutations belong to
the two pure symmetric subgroups, so the generated
group contains the single crossing transposition.
Together these groups generate $\Sym(W)$.

We now analyze the first omission, rather than
assuming its form. On $A$ reverse the left indices:
$t_i=\rho_{j-1-i}|_A$ for $0\le i<h$.
The exhausted path is $a_0,a_1,\ldots,a_h$.
Its complete proof gives
\[
 \gen{t_i,\ldots,t_{h-1}}
       =\Sym(\{a_i,\ldots,a_h\}).
\]
It also determines the individual marks: each
$t_i$ is its path transposition, with only one
possible extra transposition. For $i=h-3$ it
may additionally exchange $a_{h-1},a_h$, when
$h\ge3$. Indeed the only possible extra action
is on $\{a_{i+2},\ldots,a_h\}$ and centralizes
its full symmetric group; that centralizer is
nontrivial only when the set has size two.

The first omitted mark is $\rho_m=t_{h-1}$.
If the extra transposition just described were
present, the remaining left marks would still
have a connected orbit graph on $A$: their
ordinary path edges join $a_0$ through $a_{h-1}$,
and the extra edge joins $a_{h-1}$ to $a_h$.
With the crossing and the full right action this
would make $Q_{>m}$ transitive on $W$, contrary
to its hypothesis. Hence no extra is present,
and the remaining left marks give the pure
$\Sym(A\setminus\{a_h\})$ while fixing $a_h$.

For $h\ge3$, $\rho_j$ centralizes the pure
$\Sym(\{a_1,\ldots,a_h\})=Q_{\le j-2}$,
so fixes its entire support. Here the equality
refers to the left marks $m,\ldots,j-2$ inside
$Q$. For $h=2$, its only possible additional
action on $A$ is $(a_1\ a_2)$. If that were
present, the surviving left edge $(a_0\ a_1)$,
the crossing and this extra edge would likewise
make $Q_{>m}$ transitive. Thus it is absent.
In every case $\rho_j$ fixes $a_h$.

All surviving generators therefore fix $z=a_h$.
They contain the pure symmetric groups on
$A\setminus\{z\}$ and $B$. Multiplying
$\rho_j$ by its internal $B$-restriction gives
the single crossing transposition, since its
only action on $A$ is the crossing. These
generate the full symmetric group on
$W\setminus\{z\}$, proving the final equality.
\end{proof}

\paragraph{Attribution and scope.}
The complete anchored-path theorem, its exhausted
symmetric images, and the local crossing orbit bound
are acquired. The short-action centralizer argument
and the exact first-omission analysis are the new
steps. This is a general equality theorem with both
sides of length at least two. It makes no claim
about the remaining unequal-end separator, the
short side cases, or actual attachments by outside
generators. No original, question, numerical or
formalization credit is claimed.
\end{document}
